\chapter{Bilinear games (Section 6, Appendices E and F)}
\label{chap:bilinear}

\section{Uninformed learning with Tsallis-FTRL-SPM}

\begin{lemma}[PSMR of Tsallis-FTRL-SPM]
  \label{lem:tsallis_spm_psmr}
  \lean{Ito2026Adversarial.psmr_tsallisSPMPaper_le_sub,
    Ito2026Adversarial.psmr_tsallisSPMPaper_le_of_isStrictPSNE, Ito2026Adversarial.le_of_self_bound,
    Ito2026Adversarial.isSPMParams_paper, Ito2026Adversarial.rpow_one_sub_spmAlpha_le}
  \leanok
  \uses{lem:tsallis_spm_regret, lem:regret_relations, lem:sqrt_func, lem:self_bound}
  For a bilinear game with $1 \le d_x \le m_x$, $m_x \ge 2$, and the parameters of
  Theorem~\ref{thm:tsallis_bilinear}, with $L = \log m_x$: $\PSMR_T \le 8(16 + 8c)\sqrt{T d_x L} +
  (128c + \frac{64}{c} + \frac{16 + 8c}{c}) m_x L^2 - \Dmix T$, and for a strict PSNE $(x^\star,
  y^\star)$, $\PSMR_T \le (32(16 + 8c))^2 \frac{d_x L (1 + \log T)(1 + 1/\DCmin)}{\DRmin} + 2(128c
  + \frac{64}{c} + (16 + 8c)(48 + \frac1c)) m_x L^2$.
\end{lemma}

\begin{proof}
  \leanok
  \uses{lem:tsallis_spm_regret, lem:regret_relations, lem:sqrt_func, lem:self_bound}
  The first bound from $\PSMR_T = \mathrm{NR}_T - \Dmix T \le \ER_T - \Dmix T$ and
  Lemma~\ref{lem:tsallis_spm_regret}. For the second, with $\Delta = \DR$ and $C =
  \E[\#\{t : y_t \ne y^\star\}]$: $\E[\sum_t \Delta(x_t)] \le \ER_T(x^\star) + 4C$ and $\DCmin C \le
  \ER_T(x^\star) - \PSMR_T$ (Lemma~\ref{lem:regret_relations}); if $A \le a\sqrt{Q(A + 4C)} + K$
  then $A \le a^2 Q + \sqrt{4a^2QC} + 2K$ (Lemma~\ref{lem:self_bound}) and $\sqrt{4a^2QC} - \DCmin C
  \le a^2 Q/\DCmin$ (Lemma~\ref{lem:sqrt_func}; $C = 0$ when $\DCmin = 0$).
\end{proof}

\begin{theorem}[Theorem 5]
  \label{thm:tsallis_bilinear}
  \lean{Ito2026Adversarial.forall_exists_psmr_tsallisSPMPaper_le}
  \leanok
  \uses{def:tsallis_spm, lem:tsallis_spm_regret, lem:regret_relations, lem:sqrt_func,
    lem:self_bound}
  For every $c > 0$ there is $C$ such that for every bilinear game with $m_x \ge 2$, every exploration
  distribution $p_0$ with variance ratio $c$, every adaptive adversary and every reward noise with
  conditional mean $u$ and rewards in $[-1, 1]$, Tsallis-FTRL-SPM with the parameters of the
  theorem has $\PSMR_T \le C(\sqrt{T d_x \log m_x} + m_x \log^2 m_x)$; if $(x^\star, y^\star)$ is a
  strict PSNE, $\PSMR_T \le C(d_x \log m_x (1 + \log T)(1 + \frac{1}{\DCmin})/\DRmin + m_x \log^2 m_x)$;
  and if $\Dmix > 0$, $\PSMR_T \le C(\frac{d_x \log m_x}{\Dmix} + m_x \log^2 m_x)$.
\end{theorem}

\begin{proof}
  \leanok
  \uses{lem:tsallis_spm_regret, lem:regret_relations, lem:sqrt_func, lem:self_bound,
    lem:tsallis_spm_psmr}
  As for Theorem~\ref{thm:tsallis}, with the regret bound of Tsallis-FTRL-SPM
  (Lemma~\ref{lem:tsallis_spm_regret}) in place of those of Tsallis-INF, and
  $\ER_T(x^\star) \ge \E[\sum_t \langle p_t, \DR\rangle] - 2C$ for bilinear games (features and
  $A$ of norm at most $1$). With $\alpha = 1 - \frac{1}{4 \log m_x}$,
  $\frac{m_x^{1 - \alpha}}{\alpha(1 - \alpha)} = O(\log m_x)$, and the additive term of
  Lemma~\ref{lem:tsallis_spm_regret} is $O(m_x \log^2 m_x)$ (the constant depends on $c$).
\end{proof}

\begin{remark}
  \label{rem:tsallis_bilinear_statement}
  As in Theorem~\ref{thm:tsallis}: the paper's strict-PSNE bound
  $O(\frac{d_x \log m_x}{\DRmin \DCmin} \log T + m_x \log m_x)$ follows from the stated form when
  $m_y \ge 2$, $T \ge 2$, and is false in Lean for $m_y = 1$; the no-PSNE bound is stated under
  $\Dmix > 0$ (Remark~\ref{rem:no_psne_gap}). The paper's ``suitable choice of $p_0$'' is any
  exploration distribution with variance ratio $c$, whose design matrix is positive definite
  (implicit in the paper).

  The paper's additive term $m_x \log m_x$ is false for the parameters of the theorem: the
  penalty term $\bar\beta h_0$ of \cite[Eq.~(24)]{ito2024adaptive}, with $h_0 = -\bar\psi(q_1) \le
  \frac{m_x^\alpha}{1 - \alpha}$, is $\bar\beta h_0 \le \frac{64 m_x^\alpha \log^2 m_x}{c}$, and
  this order is attained. For $d_x = 1$, $\phi(x^\star) = 1$, $\phi(x) = -1$ for the other
  $m_x - 1$ actions, $A = 1$, a single column, deterministic rewards and $p_0$ uniform, $S(p) = 1$
  for every $p$, so $g_t = \phi$ and the run is deterministic. The barrier $\bar\beta
  \varphi_{1 - \alpha}$ keeps $\hat p_t(x) \ge e^{-1} \bar\beta / (4t + O(\log m_x))$ on each bad
  action while $\beta_t = O(t)$, so that $\PSMR_{m_x^2}$ is of order $m_x \log^2 m_x$ (simulation:
  the ratio to $\sqrt{T \log m_x} + m_x \log m_x$ at $T = m_x^2$ is $19$ for $m_x = 100$ and $37$
  for $m_x = 1000$). The bounds are stated with $m_x \log^2 m_x$.
\end{remark}

\section{Informed learning with Maximin-LinUCB}

\begin{lemma}[Confidence ellipsoid, Lemma 13]
  \label{lem:conf_ellipsoid}
  \lean{Ito2026Adversarial.probReal_forall_sqrt_mahalanobisSq_le_ge}
  \leanok
  \uses{def:maximin_linucb, lem:self_normalized, lem:log_det_le}
  In a run of any informed player in a bilinear game, against any adaptive adversary, with reward
  noise of conditional mean $u$ and rewards in $[-1, 1]$, for $\lambda > 0$ and $\delta \in (0, 1)$,
  with probability at least $1 - \delta$, for all $t$, $\norm{\vect(A) - V_t^{-1} b_t}_{V_t} \le
  \beta_t$.
\end{lemma}

\begin{proof}
  \leanok
  \uses{lem:self_normalized, lem:log_det_le, lem:bounded_noise_subgaussian, lem:ridge_error}
  $b_t = (V_t - \lambda I)\vect(A) + S_t$ with $S_t = \sum_{s < t} \eps_s a_s$, so
  $\vect(A) - V_t^{-1} b_t = V_t^{-1}(\lambda \vect(A) - S_t)$ and
  $\norm{\cdot}_{V_t} \le \norm{S_t}_{V_t^{-1}} + \sqrt\lambda \norm{\vect(A)}$
  (Lemma~\ref{lem:ridge_error}, with $u(x, y) = \langle \vect(x y^\top), \vect(A)\rangle$). The
  noise is conditionally $1$-sub-Gaussian (Lemma~\ref{lem:bounded_noise_subgaussian});
  Lemma~\ref{lem:self_normalized} bounds
  $\norm{S_t}_{V_t^{-1}}^2 \le 2 \log(1/\delta) + \log(\det V_t/\lambda^d)$, and
  Lemma~\ref{lem:log_det_le} bounds the last term by $d \log(1 + \frac{t}{d\lambda})$ ($\norm{a_s}
  = \norm{x_s}\norm{y_s} \le 1$). Finally $\norm{\vect(A)}^2 = \norm{A}_F^2 \le d_y \le d$: the
  columns $A e_j$ have norm at most $\norm{A}_2 \le 1$.
\end{proof}

\begin{lemma}[Optimism of the linear UCB index]
  \label{lem:linucb_optimism}
  \lean{Bandits.Linear.abs_inner_sub_inner_ridgeEstimate_le, Bandits.Linear.inner_le_ucbIndex,
    Bandits.Linear.ucbIndex_le_inner_add}
  \leanok
  If $V$ is positive definite, $\hat\theta = V^{-1} b$ and $\norm{\theta - \hat\theta}_V \le \beta \le
  \beta'$, then for every $a$, $\abs{\langle a, \theta\rangle - \langle a, \hat\theta\rangle} \le \beta
  \norm{a}_{V^{-1}}$, so that the index $U = \langle a, \hat\theta\rangle + \beta' \norm{a}_{V^{-1}}$
  satisfies $\langle a, \theta\rangle \le U \le \langle a, \theta\rangle + 2 \beta' \norm{a}_{V^{-1}}$.
\end{lemma}

\begin{proof}
  \leanok
  Cauchy--Schwarz for the inner product of $V$:
  $\langle a, \theta - \hat\theta\rangle^2 \le \norm{a}^2_{V^{-1}} \norm{\theta - \hat\theta}^2_V$.
\end{proof}

\begin{lemma}[One round of Maximin-LinUCB]
  \label{lem:purelinucb_round}
  \lean{Ito2026Adversarial.pureMaximin_sub_le_of_mem_linGoodEvent}
  \leanok
  \uses{def:maximin_linucb, lem:conf_ellipsoid, lem:linucb_optimism, lem:maximin_optimism}
  On the event of Lemma~\ref{lem:conf_ellipsoid} ($\lambda = 1$, $\delta = 1/T$), in the run of
  Maximin-LinUCB with $\lambda = 1$ and the radii of $\delta = 1/T$, with $d = d_x d_y \ge 1$, for
  every round $t < T$ (counted from $0$, played with the state $V_t$ of the first $t$ rounds and the
  radius $\beta_{t+1}$), $\vstar - u(x_t, y_t) \le 2 \beta_T \sqrt{\min(1, \norm{a_t}^2_{V_t^{-1}})}$.
\end{lemma}

\begin{proof}
  \leanok
  \uses{lem:linucb_optimism, lem:maximin_optimism}
  $u(x, y) = \langle \vect(x y^\top), \vect(A)\rangle$, $\norm{\vect(A) - \hat\theta_t}_{V_t} \le
  \beta_t \le \beta_{t+1}$ ($\beta$ is nondecreasing), so by Lemma~\ref{lem:linucb_optimism} $u \le
  U_t$ and Lemma~\ref{lem:maximin_optimism} gives $\vstar \le U_t(x_t, y_t) \le u(x_t, y_t) + 2
  \beta_{t+1} \norm{a_t}_{V_t^{-1}}$, with $\beta_{t+1} \le \beta_T$. Moreover $\vstar - u \le 2 \le
  2 \beta_T$ ($\beta_T \ge \sqrt d \ge 1$), whence the minimum with $1$.
\end{proof}

\begin{lemma}[Pathwise regret of Maximin-LinUCB]
  \label{lem:purelinucb_pathwise}
  \lean{Ito2026Adversarial.sum_pureMaximin_sub_le_of_mem_linGoodEvent,
    Ito2026Adversarial.sum_pureMaximin_sub_le_div_of_mem_linGoodEvent}
  \leanok
  \uses{lem:purelinucb_round, lem:elliptical_potential}
  On the same event, for $T \ge 2$, $d = d_x d_y \ge 1$:
  $\sum_{t < T} (\vstar - u(x_t, y_t)) \le 2 \beta_T \sqrt{T} \sqrt{4 d \log T}$ and, for every gap
  threshold $\Dlin > 0$, $\sum_{t < T} (\vstar - u(x_t, y_t)) \le 4 \beta_T^2 \cdot 4 d \log T /
  \Dlin$.
\end{lemma}

\begin{proof}
  \leanok
  \uses{lem:purelinucb_round, lem:elliptical_potential}
  Lemma~\ref{lem:elliptical_potential} ($\lambda = 1$, $\norm{a_t} = \norm{x_t} \norm{y_t} \le 1$)
  bounds $\sum_{t < T} \min(1, \norm{a_t}^2_{V_t^{-1}})$ by $2 d \log(1 + T/d) \le 4 d \log T$.
  The worst-case bound follows from Lemma~\ref{lem:purelinucb_round} by Cauchy--Schwarz, the
  instance-dependent one from $\vstar - u \le (\vstar - u)^2/\Dlin$ (as $\vstar - u \notin (0,
  \Dlin)$).
\end{proof}

\begin{theorem}[Theorem 12]
  \label{thm:purelinucb}
  \lean{Ito2026Adversarial.exists_psmr_maximinLinUCB_le}
  \leanok
  \uses{def:maximin_linucb, lem:conf_ellipsoid, lem:purelinucb_pathwise, def:regrets}
  There is a universal constant $C$ such that for every bilinear game, every $T \ge 2$, every
  adaptive adversary and every reward noise with conditional mean $u$ and rewards in $[-1, 1]$,
  Maximin-LinUCB with $\lambda = 1$ and the radii $\beta_t$ of $\delta = 1/T$ has
  $\PSMR_T \le C d_x d_y \sqrt T \log T$ and $\PSMR_T \le C \frac{d_x^2 d_y^2 \log^2 T}{\Dlin}$ for
  every gap threshold $\Dlin > 0$.
\end{theorem}

\begin{proof}
  \leanok
  \uses{lem:conf_ellipsoid, lem:purelinucb_pathwise}
  If $d = d_x d_y = 0$ the game is zero and $\PSMR_T = 0$. Otherwise, with $L = \log T \ge 1/2$,
  $\beta_T^2 \le 12 d L$, and Lemma~\ref{lem:purelinucb_pathwise} bounds the regret on the event of
  Lemma~\ref{lem:conf_ellipsoid} by $14 d \sqrt T L$ and $192 d^2 L^2/\Dlin$. The failure event, of
  probability at most $1/T$, contributes at most $2$ (the regret is at most $2T$), respectively
  $4/\Dlin$ (the regret is at most $4T/\Dlin$, since $\vstar - u \le (\vstar - u)^2/\Dlin \le
  4/\Dlin$). This gives $C = 256$.
\end{proof}

\begin{theorem}[Theorem 6]
  \label{thm:purelinucb_informal}
  \lean{Ito2026Adversarial.exists_forall_exists_psmr_maximinLinUCB_le}
  \leanok
  \uses{thm:purelinucb}
  There is a universal constant $C$ such that for every bilinear game and every $T \ge 2$ there are
  parameters $(\lambda, \beta)$ of Maximin-LinUCB with $\PSMR_T \le C d_x d_y \sqrt T \log T$ and
  $\PSMR_T \le C \frac{d_x^2 d_y^2 \log^2 T}{\Dlin}$, against every adaptive adversary and reward
  noise.
\end{theorem}

\begin{proof}
  \leanok
  \uses{thm:purelinucb}
  Theorem~\ref{thm:purelinucb}, with $\lambda = 1$ and the radii of Theorem~\ref{thm:purelinucb}.
\end{proof}
